\scanpage{155}
\deleted{Décomposition canonique}

\uline{Niveau d'un motif}

\unsure{a)} \uline{Motif effectif} $M$ de degré $i\geq0$ est dit de niveau $\leq i-k\rho$ s'il est isomorphe à un motif de la forme $M'(\unsure{k|-k})$, où $M'$ est un motif effectif de degré $i-k\rho$.
Les motifs effectifs, de degré $i\geq0$ fixé, de niveau $\leq i-k\rho$ font une sous-catégorie \uline{épaisse} de $\mathcal M^i$, donc tout objet\inserted{ $M$} de $\mathcal M^i$ se décompose canoniquement en une somme d'un\inserted{ tel} objet maximal $M'$ de niveau $\leq i-k\rho$, et d'un \inserted{autre sous-}objet. \deleted{\unsure{M}}
Ainsi,
\[
\mathcal M^i\simeq\mathcal P^i\times\mathcal P^{i-\rho}(\unsure{-1|1})\times\mathcal P^{i-2\rho}(\unsure{-2|2})\times\cdots,
\]
\margin{projections $\kappa_{i,k}$}
où $\mathcal P^j$ désigne la catégorie des motifs effectifs de degré $j$ \deleted{primitifs} dont aucun sous-motif\inserted{ $\ne0$} n'est de niveau $\leq j-\rho$.
\margin{motifs \uline{primitifs de niveau $j$}}

Celà donne donc une décomposition de $\mathcal M^{*}=\mathcal M$ en termes des $\mathcal P^i$ (pour les structures additives). La structure \unsure{multiplication|multiplicative} est connue quand on connaît les foncteurs
\diagram{\begin{tikzcd}
\unsure{\theta|\Theta}^{i,j,k}:\mathcal P^i\times\mathcal P^j\arrow[r] & \mathcal P^{i+j-k\rho}
\end{tikzcd}}
\margin{$\unsure{\theta|\Theta}^{ijk}(M,N)=\deleted{\unsure{\kappa(M\otimes N(-}}\kappa_{i,k}(M\otimes N)(k)$.}

b) Niveau d'un motif quelconque (pas néc. effectif).
Se définit en se ramenant au cas précédent....
\mspage{154}
